% FORMALIZACION_REUNIDA of the calendar projection and twelve-position
% representation already developed in the TPK architecture and the Witt link.
% Explicit formalization of the Gram comparison and frame transport.
\subsubsection{Calendar projection and representation of its twelve positions}
\label{ii:dodeca:calendario}
The calendar of Section~\ref{ii:tpk-arquitectura} provides an
observable coordinate of the complete state. In its notation, write
\(r_0(\theta)=r(\theta)-1\in\{0,\ldots,8\}\) and
\(b=\beta(\theta)\in\mathbb Z/12\mathbb Z\). The clock projection is
\[
 Q_{\rm clk}(x)=\bigl(\beta(\theta(x)),r_0(\theta(x))\bigr).
\]
For an admissible route of \(n\) updates, the induced transport is
\begin{equation}
 k(r_0,n)=\left\lfloor\frac{r_0+n}{9}\right\rfloor,
 \qquad
 T_n(b,r_0)=\bigl(b+k(r_0,n)\bmod12,\ (r_0+n)\bmod9\bigr).
 \label{ii:dodeca:transporte-reloj}
\end{equation}
Here \(n\in\mathbb N_0\) counts effective updates; the integer
quotient \(k\) records the nonadic returns crossed.

\begin{proposition}[Compatibility of the calendar with route composition]
\label{ii:dodeca:composicion}
The projection \(Q_{\rm clk}\), together with the length of each route, defines
a functor from the category of admissible TPK routes to the action category
of calendar~\eqref{ii:dodeca:transporte-reloj}. The identities
\(T_{m+n}=T_mT_n\), \(T_0=I\),
\(T_9(b,r_0)=(b+1,r_0)\), and \(T_{108}(b,r_0)=(b,r_0)\) hold.
\end{proposition}
\begin{proof}
Putting \(r'_0=(r_0+n)\bmod9\), Euclidean division gives
\[
 k(r_0,n+m)=k(r_0,n)+k(r'_0,m).
\]
This identity simultaneously preserves residue and quotient when the
routes are concatenated. The update of \(\theta\) established in
Section~\ref{ii:tpk-arquitectura} induces exactly
\eqref{ii:dodeca:transporte-reloj}; the remaining identities follow
by substitution. The target category has pairs \((b,r_0)\) as
objects and admissible lengths as arrows, with addition as composition.
\end{proof}

The index \(p=b+1\in\{1,\ldots,12\}\), taking the representative
\(0\le b<12\), indicates the dodecaphase position. One turn of
\(108\) updates traverses its twelve positions. The proved periodicity
belongs to this projection: route, accumulated carry, and
memory remain in the source state and continue to be updated.

\paragraph{Incidence realization from the angular projector.}
The representatives \(r_pC_5\) and representation \(\rho\) of
Section~\ref{sec:ii-enlace-angular-witt} fix the chart of the twelve
positions. The projector already defined by
\eqref{eq:ii-witt-projector-character} produces
\begin{equation}
 \begin{gathered}
 \mathcal H_3=P_3^{\mathrm{ico}}\mathbb R^{12},\qquad v_p=2P_3^{\mathrm{ico}}e_p,
 \\[-2pt]
 (P_3^{\mathrm{ico}})_{pq}=\frac1{20}
 \sum_{\substack{g\in A_5\\gr_qC_5=r_pC_5}}\chi_3(g^{-1}).
 \end{gathered}
 \label{ii:dodeca:vertices-proyector}
\end{equation}
The last sum contains five terms per entry and uses the
representatives and the two classes of order-five cycles fixed there.

\begin{proposition}[Icosahedral realization of the dodecaphase positions]
\label{ii:dodeca:realizacion}
The vectors \(v_p\) are twelve distinct unit vectors in
\(\mathcal H_3\). Their incidence is
\(p\sim q\Longleftrightarrow\langle v_p,v_q\rangle=1/\sqrt5\).
The map \(p\mapsto v_p\) is \(A_5\)-equivariant, and its image
is isometric to \(\Omega_{12}^{\rm ico}\).
\end{proposition}
\begin{proof}
Transitivity and \(\operatorname{Tr}P_3^{\mathrm{ico}}=3\) give
\((P_3^{\mathrm{ico}})_{pp}=1/4\); by idempotence,
\(\langle v_p,v_q\rangle=4(P_3^{\mathrm{ico}})_{pq}\) and \(\|v_p\|=1\).
Evaluation of the five terms in
\eqref{ii:dodeca:vertices-proyector} gives the antipodal pairs
\[
 (1,4),\ (2,3),\ (5,8),\ (6,7),\ (9,12),\ (10,11).
\]
For the representatives \((1,2,5,6,9,10)\), their Gram matrix is
\begin{equation}
 I_6+\frac D{\sqrt5},\qquad
 D=\begin{pmatrix}
 0&-1&-1&1&-1&1\\
 -1&0&1&1&-1&-1\\
 -1&1&0&1&1&1\\
 1&1&1&0&-1&1\\
 -1&-1&1&-1&0&1\\
 1&-1&1&1&1&0
 \end{pmatrix},\qquad D^2=5I_6.
 \label{ii:dodeca:gram}
\end{equation}
The remaining inner products are obtained by antipodality.
Thus, apart from a vertex and its opposite, the products are
\(\pm1/\sqrt5\), and each vertex has five neighbors. The following table
makes explicit the comparison with the signed six-position chart
\(\mathcal U=(\infty,0,1,2,3,4)\):
\begin{equation}
\begin{array}{c|cc@{\qquad}c|cc}
p&u(p)&\sigma(p)&p&u(p)&\sigma(p)\\ \hline
1&\infty&+1&7&3&-1\\
2&0&-1&8&1&+1\\
3&0&+1&9&2&-1\\
4&\infty&-1&10&4&+1\\
5&1&-1&11&4&-1\\
6&3&+1&12&2&+1
\end{array}
\label{ii:dodeca:carta-firmada}
\end{equation}
Indeed, for \(C=\operatorname{bal}(A_W)\), with \(A_W\) given in
\eqref{exc:aw}, the entries of \(D\) are
\(D_{ab}=\sigma(p_a)\sigma(p_b)C_{u(p_a),u(p_b)}\).
Consequently, the table and~\eqref{ii:dodeca:gram} isometrically
identify \(v_p\) with
\(\sigma(p)\sqrt2 E_Ce_{u(p)}\), where
\(E_C=(I_6+C/\sqrt5)/2\). Both systems span three-dimensional
spaces and have identical Gram matrices, proving that the induced
linear map is well defined and isometric. The explicit
realization in~\eqref{ii:ico:proyector-seis}, proved below,
identifies it with~\(\Omega_{12}^{\rm ico}\) and its incidence.
Finally, \(P_3^{\mathrm{ico}}\rho(g)=\rho(g)P_3^{\mathrm{ico}}\) implies
\(v_{g\cdot p}=\rho(g)v_p\). The bijection of the twelve positions gives
\(\operatorname{Stab}_{A_5}(v_p)=r_pC_5r_p^{-1}\).
\end{proof}

One thus obtains the representational composition
\begin{equation}
 x\longmapsto Q_{\rm clk}(x)\longmapsto
 p=\beta(\theta(x))+1\longmapsto v_p.
 \label{ii:dodeca:composicion-representacional}
\end{equation}
The signed expression of this composition is
\(r_\pm^{\rm rep}(x)=(u(p),\sigma(p))\), determined by
Table~\eqref{ii:dodeca:carta-firmada}. The table is an explicit choice
of alignment between charts; the actions of \(A_5\) provide its
equivalent alignments. The complete record is preserved in the
domain of~\eqref{ii:dodeca:composicion-representacional}.

\paragraph{Link with the five-component tensor carrier.}
In \(\mathcal H_3\), form
\(Q_p=v_pv_p^*-I_{\mathcal H_3}/3\). Let \(J=\mathbf1\mathbf1^*\),
and let \(\mathsf A\) be the matrix of the preceding antipodal permutation.
The Frobenius product and~\eqref{ii:dodeca:gram} give
\begin{equation}
 \begin{aligned}
 \operatorname{Gram}(Q_p)
 &=(4P_3^{\mathrm{ico}})\circ(4P_3^{\mathrm{ico}})-\frac13J
 =\frac45(I+\mathsf A)-\frac2{15}J,
 \\
 P^{\rm ico}_5&=\frac58\operatorname{Gram}(Q_p)
 =\frac12(I+\mathsf A)-\frac1{12}J.
 \end{aligned}
 \label{ii:dodeca:puente-tensorial}
\end{equation}
The symbol \(\circ\) denotes the entrywise product. The last
expression is the orthogonal projector onto coordinates even under
antipodality and orthogonal to \(\mathbf1\), of dimension five. Its character
is \((5,1,-1,0,0)\): on the six antipodal pairs, the fixed-point
counts are \((6,2,0,1,1)\), and the constant component is removed.
It therefore coincides with the central projector of
\eqref{ii:pent:eq:gravity-p5}. This identity links the three-dimensional
construction with \(\Gamma_0\) and \(\Gamma_5\), defined in
\eqref{ii:pent:eq:gravity-Gamma0} and~\eqref{ii:pent:eq:gravity-gamma5},
through the explicit quadratic operation \(v_p\mapsto Q_p\).

\subsubsection{Transport of the representational frame}
\label{ii:dodeca:marco}
Let \(s=(1\,2\,\cdots\,12)\), and let \(S_ce_p=e_{s(p)}\).
A nonadic calendar return takes \(p\) to \(s(p)\). The joint
transport of frame, representation, and projector is expressed by
\begin{equation}
 \begin{gathered}
 P_{3,k}^{\mathrm{ico}}=S_c^kP_3^{\mathrm{ico}}S_c^{-k},\qquad
 \rho^{(k)}(g)=S_c^k\rho(g)S_c^{-k},\\
 v^{(k)}_{s^k(p)}=2P_{3,k}^{\mathrm{ico}}e_{s^k(p)}=S_c^kv_p.
 \end{gathered}
 \label{ii:dodeca:marco-movil}
\end{equation}
The last equality follows by substitution; orthogonality of
\(S_c\) preserves all inner products and transports
incidence. For a route of length \(n\), take
\(k=k(r_0,n)\). The carry identity in
Proposition~\ref{ii:dodeca:composicion} proves compatibility of
these transports with concatenation. Although \(S_c^{12}=I\),
the twelve turns update the memory of the source route.

Equality~\eqref{ii:dodeca:marco-movil} establishes covariance between
frames. Active action on a fixed incidence has additional content:
it requires identifying transport of the complete state and its
action in that chart. In particular, \(s\) has order twelve, whereas
the orders of elements of \(A_5\) are \(1,2,3,5\).
Likewise, the stabilizer of a vertex has order five, and that of its
antipodal pair has order ten; the involution of the signed chart
exchanges the two vertices. These distinctions preserve the meaning
of orientation, sheet, and return in properties (i)--(iv): the
proved representational realization and the descent of the complete exceptional
route are related through those properties, with their
domains and fibers declared.
